\scanpage{193}
% manuscript position: top
\mspage{109}
% manuscript position: 1–3
\deleted{des cinq conditions, autres que $c$, soit $P(\pi)$ la condition obtenue et l'implication $P(\pi)\Rightarrow c$. Il restera à prouver encore}
% manuscript position: 4
\[
\deleted{c\Rightarrow P(\pi).}
\]

% manuscript position: 5
\uline{Prouvons} $a''\Rightarrow a$. Considérons l'application
% manuscript position: 6
\[
\operatorname{Fl} C_{/F}\longrightarrow\operatorname{Fl} C\tag*{(\unsure{\(\alpha\)|\(\ast\)})}
\]
% manuscript position: 7–16
induite par le foncteur \inserted{étale} $C'=C_{/F}\longrightarrow C$, alors l'ens. des flèches au dessus d'une flèche $\alpha:x\longrightarrow y$ de $C$ est \deleted{de cardinal} en corr. biunivoque avec $\operatorname{Ob} C'_y=F(y)$, puisqu'une telle flèche $\alpha'$ est connue quand on connaît son but, qui est un objet quelconque de $F(y)$ ($C'\longrightarrow C$ étant \unsure{étale} i.e. fibré à fibres discrètes). Donc la fibre de \unsure{(\(\alpha\))|(\(\ast\))} en $\alpha$ est \deleted{majorée} de cardinal \deleted{$\leq$} $\operatorname{card} F(y)\leq\pi$, donc on a
% manuscript position: 17
\[
\operatorname{Card}\operatorname{Fl} C_{/F}\leq\pi\underbrace{\operatorname{card}\operatorname{Fl} C}_{\leq\pi_1}
\leq\pi\underbrace{\pi_1}_{\leq\pi}\leq\pi^2=\pi,
\]
% manuscript position: 18
ce qui prouve $a$.

% manuscript position: 19–21
Prouvons enfin $c\Rightarrow a$. Pour ceci, nous allons utiliser \inserted{la première assertion \unsure{du n°|du no} $M_\pi\subset\operatorname{Filt}^{\unsure{\geq\pi|\pi}} M$ de} SGA 4 I 9.15, \deleted{qui} s'applique dès lors que nous prouvons le
% manuscript position: left/22–27
% slanted
\margin{\unsure{2|Z} Pas prouvé ; la condition c) \unsure{aux filtrations|aux filtrations sur} \unsure{cardinales|cardinaux} paraît-être \unsure{pas satisfaite|en défaut}. \deleted{\unsure{Sans doute}} Faut \unsure{affiner|rectifier} \deleted{\illegible{1}}}

% manuscript position: 22–27
\uline{Corollaire} 4.5.3. Sous les conditions de 4.5.2, soit \deleted{\unsure{,}} pour tout cardinal $\pi\geq\pi_1$, $\deleted{\unsure{M_\pi}}\operatorname{Filt}^{\pi}M$ la sous-\uline{\deleted{catégorie strictement pleine}} \inserted{\deleted{\unsure{ensemble}}} de $\operatorname{Ob} M$ formé des $F\in\operatorname{Ob}M$ \deleted{\unsure{$\pi$-accessibles (une des conditions)}} qui satisfont la condition a) de 4.5.2, Alors la famille des $(\operatorname{Filt}^{\pi}M)_{\pi\geq\pi_1}$ est une
% manuscript position: foot
% archive sheet number: 192
